% Internal obstruction record. This is not a claim that odd addition is false.
\subsection*{Why the existing odd-plus-even boundary does not close the image}
Let $a$ be odd, $b$ even, $d=a+b$, and $f\in\{0,1\}$.
The target odd relative rank is $R=9d^2-1$. A transverse union of
sharp components in degrees $a,b$ has relative stable rank
\[
 (9a^2-1)+(9b^2-3+f)+f+(18ab-1)
       =R-(4-2f).
\]
Thus this union does not lie in the maximal-rank target open.
The available natural odd block, the common kernel of the even
component's invariant covectors, and the intersection heart have
total dimension
\[
 (9a^2-1)+(9b^2-4)+(18ab-2)=R-6.
\]
Even if all their gluing translations produce an actual natural
$\GL_{R-6}(\mathbf F_2)$ block, this does not prove irreducibility
of the full $R$-dimensional representation.

The concrete countermodel communicated by the reduction/descent
reviewer is the proper parabolic
\[
 P=\Hom(\mathbf F_2^6,\mathbf F_2^{R-6})
       \rtimes\bigl(\GL_{R-6}(\mathbf F_2)
                      \times\GL_6(\mathbf F_2)\bigr).
\]
It preserves the proper $(R-6)$-space. It has no invariant vectors
and no invariant covectors: the diagonal general linear factors
already have neither on their natural representations or duals.
It contains the actual large natural block, all rank-one columns
into that block, and the previously obtained smaller blocks.
The boundary subgroup may act trivially on a chosen
$(2+2f)$-dimensional subspace of its six-dimensional quotient,
which is the remaining boundary rank above the $(R-6)$-space.
Consequently none of these subgroup or fixed-space checks excludes
this proper image. For $a=3,b=2,f=0$, it is the explicit
$\Hom(\mathbf F_2^6,\mathbf F_2^{218})\rtimes
(\GL_{218}\times\GL_6)$ subgroup of $\GL_{224}$.

This records a limitation of the available proof, not a geometric
counterexample to full monodromy. A new argument must control the
residual quotient, or supply boundary comparisons that rule out the
proper large subspace. Clopen Lang closure alone does not preserve
the full dimension of an arbitrary invariant space across this
rank drop. The earlier parent summary is retained in
\path{notes/ROOT_SEARCH_LOG_24.md}, under the odd-plus-even paragraph.
